% Figura: hmi-alarmas -- requiere % =====================================================================
%  Estilos compartidos por los diagramas TikZ de los capitulos.
%  Va en el PREAMBULO del documento principal:
%      % =====================================================================
%  Estilos compartidos por los diagramas TikZ de los capitulos.
%  Va en el PREAMBULO del documento principal:
%      % =====================================================================
%  Estilos compartidos por los diagramas TikZ de los capitulos.
%  Va en el PREAMBULO del documento principal:
%      \input{diagramas/estilos-diagramas}
% =====================================================================

\usepackage{tikz}
\usepackage{amsmath}
\usepackage{amssymb}
\usetikzlibrary{arrows.meta,positioning,shapes.geometric,shapes.misc,calc,fit,backgrounds,decorations.pathreplacing}

\definecolor{cbloque}{HTML}{E8EEF4}
\definecolor{cborde}{HTML}{4A6D8C}
\definecolor{cacento}{HTML}{D9E6D5}
\definecolor{cacentob}{HTML}{5A7A4A}
\definecolor{cgris}{HTML}{EDEDED}
\definecolor{cgrisb}{HTML}{888888}
\definecolor{calarma}{HTML}{F6DEDE}
\definecolor{calarmab}{HTML}{A03030}

\tikzset{
  bloque/.style={draw=cborde,fill=cbloque,rounded corners=2pt,align=center,
                 inner sep=5pt,font=\footnotesize,minimum height=9mm},
  bloqueg/.style={bloque,draw=cgrisb,fill=cgris},
  bloquea/.style={bloque,draw=cacentob,fill=cacento},
  bloquer/.style={bloque,draw=calarmab,fill=calarma},
  proceso/.style={bloque,minimum width=52mm},
  decision/.style={draw=cborde,fill=cbloque,align=center,font=\scriptsize,
                   diamond,aspect=2.2,inner sep=1pt},
  flecha/.style={-{Stealth[length=2mm]},draw=cborde!80!black,thick},
  flechag/.style={flecha,draw=cgrisb},
  et/.style={font=\scriptsize,inner sep=1.5pt,fill=white,align=center},
  etl/.style={et,midway},
  grupo/.style={draw=cgrisb,dashed,rounded corners=3pt,inner sep=4mm},
  tit/.style={font=\scriptsize\bfseries,text=cgrisb!60!black},
}

% --- utilidades para diagramas de secuencia ---
\tikzset{
  vida/.style={draw=cgrisb,dash pattern=on 2pt off 2pt},
  cab/.style={bloque,minimum width=26mm,font=\scriptsize\bfseries},
  nota/.style={draw=cacentob,fill=cacento,font=\scriptsize,align=center,
               rounded corners=1pt,inner sep=3pt},
}
\newcommand{\msg}[4]{\draw[flecha] (#1,#3) -- (#2,#3) node[et,midway,above=0pt]{#4};}
\newcommand{\msgr}[4]{\draw[flecha,dashed] (#1,#3) -- (#2,#3) node[et,midway,above=0pt]{#4};}
\newcommand{\auto}[3]{%
  \draw[flecha] (#1,#2) -- ++(0.55,0) -- ++(0,-0.42) -- (#1,{#2-0.42});
  \node[et,anchor=west] at ({#1+0.7},{#2-0.21}) {#3};}


% =====================================================================

\usepackage{tikz}
\usepackage{amsmath}
\usepackage{amssymb}
\usetikzlibrary{arrows.meta,positioning,shapes.geometric,shapes.misc,calc,fit,backgrounds,decorations.pathreplacing}

\definecolor{cbloque}{HTML}{E8EEF4}
\definecolor{cborde}{HTML}{4A6D8C}
\definecolor{cacento}{HTML}{D9E6D5}
\definecolor{cacentob}{HTML}{5A7A4A}
\definecolor{cgris}{HTML}{EDEDED}
\definecolor{cgrisb}{HTML}{888888}
\definecolor{calarma}{HTML}{F6DEDE}
\definecolor{calarmab}{HTML}{A03030}

\tikzset{
  bloque/.style={draw=cborde,fill=cbloque,rounded corners=2pt,align=center,
                 inner sep=5pt,font=\footnotesize,minimum height=9mm},
  bloqueg/.style={bloque,draw=cgrisb,fill=cgris},
  bloquea/.style={bloque,draw=cacentob,fill=cacento},
  bloquer/.style={bloque,draw=calarmab,fill=calarma},
  proceso/.style={bloque,minimum width=52mm},
  decision/.style={draw=cborde,fill=cbloque,align=center,font=\scriptsize,
                   diamond,aspect=2.2,inner sep=1pt},
  flecha/.style={-{Stealth[length=2mm]},draw=cborde!80!black,thick},
  flechag/.style={flecha,draw=cgrisb},
  et/.style={font=\scriptsize,inner sep=1.5pt,fill=white,align=center},
  etl/.style={et,midway},
  grupo/.style={draw=cgrisb,dashed,rounded corners=3pt,inner sep=4mm},
  tit/.style={font=\scriptsize\bfseries,text=cgrisb!60!black},
}

% --- utilidades para diagramas de secuencia ---
\tikzset{
  vida/.style={draw=cgrisb,dash pattern=on 2pt off 2pt},
  cab/.style={bloque,minimum width=26mm,font=\scriptsize\bfseries},
  nota/.style={draw=cacentob,fill=cacento,font=\scriptsize,align=center,
               rounded corners=1pt,inner sep=3pt},
}
\newcommand{\msg}[4]{\draw[flecha] (#1,#3) -- (#2,#3) node[et,midway,above=0pt]{#4};}
\newcommand{\msgr}[4]{\draw[flecha,dashed] (#1,#3) -- (#2,#3) node[et,midway,above=0pt]{#4};}
\newcommand{\auto}[3]{%
  \draw[flecha] (#1,#2) -- ++(0.55,0) -- ++(0,-0.42) -- (#1,{#2-0.42});
  \node[et,anchor=west] at ({#1+0.7},{#2-0.21}) {#3};}


% =====================================================================

\usepackage{tikz}
\usepackage{amsmath}
\usepackage{amssymb}
\usetikzlibrary{arrows.meta,positioning,shapes.geometric,shapes.misc,calc,fit,backgrounds,decorations.pathreplacing}

\definecolor{cbloque}{HTML}{E8EEF4}
\definecolor{cborde}{HTML}{4A6D8C}
\definecolor{cacento}{HTML}{D9E6D5}
\definecolor{cacentob}{HTML}{5A7A4A}
\definecolor{cgris}{HTML}{EDEDED}
\definecolor{cgrisb}{HTML}{888888}
\definecolor{calarma}{HTML}{F6DEDE}
\definecolor{calarmab}{HTML}{A03030}

\tikzset{
  bloque/.style={draw=cborde,fill=cbloque,rounded corners=2pt,align=center,
                 inner sep=5pt,font=\footnotesize,minimum height=9mm},
  bloqueg/.style={bloque,draw=cgrisb,fill=cgris},
  bloquea/.style={bloque,draw=cacentob,fill=cacento},
  bloquer/.style={bloque,draw=calarmab,fill=calarma},
  proceso/.style={bloque,minimum width=52mm},
  decision/.style={draw=cborde,fill=cbloque,align=center,font=\scriptsize,
                   diamond,aspect=2.2,inner sep=1pt},
  flecha/.style={-{Stealth[length=2mm]},draw=cborde!80!black,thick},
  flechag/.style={flecha,draw=cgrisb},
  et/.style={font=\scriptsize,inner sep=1.5pt,fill=white,align=center},
  etl/.style={et,midway},
  grupo/.style={draw=cgrisb,dashed,rounded corners=3pt,inner sep=4mm},
  tit/.style={font=\scriptsize\bfseries,text=cgrisb!60!black},
}

% --- utilidades para diagramas de secuencia ---
\tikzset{
  vida/.style={draw=cgrisb,dash pattern=on 2pt off 2pt},
  cab/.style={bloque,minimum width=26mm,font=\scriptsize\bfseries},
  nota/.style={draw=cacentob,fill=cacento,font=\scriptsize,align=center,
               rounded corners=1pt,inner sep=3pt},
}
\newcommand{\msg}[4]{\draw[flecha] (#1,#3) -- (#2,#3) node[et,midway,above=0pt]{#4};}
\newcommand{\msgr}[4]{\draw[flecha,dashed] (#1,#3) -- (#2,#3) node[et,midway,above=0pt]{#4};}
\newcommand{\auto}[3]{%
  \draw[flecha] (#1,#2) -- ++(0.55,0) -- ++(0,-0.42) -- (#1,{#2-0.42});
  \node[et,anchor=west] at ({#1+0.7},{#2-0.21}) {#3};}

 en el preambulo
\begin{tikzpicture}[x=1cm,y=1cm]

\node[bloque,text width=44mm] (est) at (0,0) {Estado del proceso\\\scriptsize (\texttt{ProcessData} tras cada ciclo)};

\node[bloquer,text width=60mm] (a1) at (7.8,2.6)
  {\textbf{ALTA} — Segundo modo detectado\\[1pt]\scriptsize \texttt{second\_mode\_alarm}: la medición no es confiable};
\node[bloquer,text width=60mm] (a2) at (7.8,1.0)
  {\textbf{ALTA} — Temperatura fuera de límites\\[1pt]\scriptsize \texttt{heater\_alarm\_priority()} $=$ \texttt{ALARM\_HIGH}};
\node[bloque,text width=60mm] (a3) at (7.8,-0.6)
  {\textbf{MEDIA} — Temperatura cerca del límite\\[1pt]\scriptsize 5\,\% final del rango};
\node[bloque,text width=60mm] (a4) at (7.8,-2.2)
  {\textbf{MEDIA} — Escrituras protegidas\\[1pt]\scriptsize \texttt{writes\_enabled is False}: usar Desbloquear (\texttt{u})};
\node[bloque,text width=60mm] (a5) at (7.8,-3.8)
  {\textbf{MEDIA} — Seeder sin respuesta\\[1pt]\scriptsize \texttt{seeder\_responding is False}};

\node[bloque,text width=40mm] (ord) at (15.8,-0.6)
  {Filtrar las activas\\y ordenar por prioridad\\[1pt]\scriptsize \texttt{ALARM\_HIGH} $=0$, \texttt{ALARM\_MEDIUM} $=1$};
\node[decision,text width=22mm] (vac) at (15.8,-3.6) {¿lista\\vacía?};
\node[bloqueg,text width=34mm] (oc) at (11.8,-5.9) {Ocultar la banda};
\node[bloquea,text width=48mm] (mo) at (18.4,-5.9)
  {Mostrar la de mayor prioridad:\\[1pt]\scriptsize prioridad + qué pasa + qué hacer,\\ con el contador ``$+n$ más''};

\foreach \n in {a1,a2,a3,a4,a5}{\draw[flechag] (est) -- (\n.west);}
\foreach \n in {a1,a2,a3,a4,a5}{\draw[flechag] (\n.east) -- (ord.west);}
\draw[flecha] (ord) -- (vac);
\draw[flecha] (vac) -| node[et,pos=0.25,above]{sí} (oc);
\draw[flecha] (vac) -| node[et,pos=0.25,above]{no} (mo);

\node[nota,text width=82mm,anchor=north west] at (0,-7.2)
 {\texttt{heater\_alarm\_priority()} es una función pura y la única fuente del criterio de temperatura:
  la usan tanto el indicador del panel como la banda de alarmas, de modo que no pueden
  discrepar. El color es redundante con el texto, nunca la única pista (ISA-18.2).};

\end{tikzpicture}
